\chapter{Iron Studies (Blood Test)}

\section*{Definition and Overview}
Iron studies are a group of blood tests that measure how much iron is in the body and how well it is being handled. The tests include serum iron, the amount of iron circulating in the blood; ferritin, which reflects the body's iron stores; transferrin, the protein that carries iron; and transferrin saturation, the percentage of transferrin that is loaded with iron. Iron studies are used to diagnose iron deficiency and iron deficiency anaemia, to distinguish it from other causes of anaemia, and to diagnose and monitor iron overload, such as haemochromatosis. They are among the most frequently requested blood tests and are usually drawn together with a full blood count.

\section*{Epidemiology}
Iron studies are ordered for many Australians every day, most commonly to investigate fatigue, anaemia, heavy periods, or pregnancy. Iron deficiency affects around one in eight women of reproductive age and is the most common cause of anaemia worldwide, while iron overload from hereditary haemochromatosis affects around one in 200 people of northern European descent, making genetic testing for it routine. The interpretation of iron studies is complicated by the fact that ferritin rises with inflammation, so results must be read alongside other markers and the clinical picture.

\section*{Aetiology and Pathophysiology}
Iron in the body exists in several compartments. Most is bound in haemoglobin in red blood cells; a smaller amount circulates in the blood bound to transferrin; and the rest is stored as ferritin, mainly in the liver, spleen, and bone marrow. In iron deficiency, stores fall first, so ferritin is low; then circulating iron falls and transferrin rises, giving a low transferrin saturation; and finally haemoglobin falls, producing anaemia. In iron overload, ferritin and transferrin saturation rise. Because ferritin is also an acute-phase protein that rises with inflammation, infection, and liver disease, it can be falsely normal in iron-deficient people who are also unwell, which is why iron studies are interpreted in context.

\section*{Clinical Features}
\begin{itemize}
    \item Iron studies are simple blood tests drawn from the arm, usually fasting or with the timing noted
    \item The key markers are serum iron, ferritin, transferrin, and transferrin saturation
    \item Low ferritin is the most specific marker of iron deficiency
    \item Low serum iron with high transferrin and low saturation supports the diagnosis
    \item High ferritin and high saturation suggest iron overload or other causes
    \item Ferritin is raised by inflammation, so a normal ferritin does not exclude deficiency when unwell
    \item Results are interpreted with the full blood count and clinical history
\end{itemize}

\section*{Differential Diagnosis}
The pattern of iron studies helps distinguish several conditions.
\begin{itemize}
    \item \textbf{Iron deficiency anaemia:} low ferritin, low iron, high transferrin, low saturation
    \item \textbf{Anaemia of chronic disease:} normal or high ferritin with low iron and low saturation
    \item \textbf{Thalassaemia:} a low mean cell volume with normal or high iron markers
    \item \textbf{B12 and folate deficiency:} anaemia with normal iron studies
    \item \textbf{Haemochromatosis:} high ferritin and high transferrin saturation
    \item \textbf{Liver disease and inflammation:} raised ferritin from causes other than iron excess
\end{itemize}

\section*{When to Seek Medical Care}
Blood tests for iron are appropriate for fatigue, paleness, breathlessness, heavy periods, pregnancy, and family history of haemochromatosis or anaemia. Results should be discussed with a doctor, because interpretation requires the full picture and treatment should not be started on a single result. Anyone with persistent symptoms or unexplained anaemia should be investigated further.

\textbf{Call 000 (or local emergency services) for:}
\begin{itemize}
    \item Severe breathlessness, chest pain, or collapse
    \item Black, tarry, or bloody stools, or vomiting blood
    \item Heavy bleeding with faintness or dizziness
    \item Sudden severe weakness or confusion
\end{itemize}

\section*{Diagnosis and Investigations}
\begin{itemize}
    \item \textbf{Iron studies:} serum iron, ferritin, transferrin, and transferrin saturation
    \item \textbf{Full blood count:} haemoglobin and red cell size, essential for interpreting iron status
    \item \textbf{Inflammatory markers:} C-reactive protein, to adjust for ferritin raised by inflammation
    \item \textbf{B12 and folate:} to exclude other deficiencies
    \item \textbf{Liver function tests:} where iron overload is suspected
    \item \textbf{Genetic testing:} for HFE gene mutations in suspected haemochromatosis
    \item \textbf{Further investigation:} endoscopy or colonoscopy for bleeding when deficiency is unexplained
\end{itemize}

\section*{Management}
\begin{itemize}
    \item \textbf{Iron deficiency:} oral iron supplements with dietary change, and investigation of the cause
    \item \textbf{Intravenous iron:} for intolerance, malabsorption, or late pregnancy
    \item \textbf{Iron overload:} regular venesection (blood removal) in haemochromatosis, and monitoring
    \item \textbf{Treatment of underlying conditions:} coeliac disease, bleeding, and chronic disease
    \item \textbf{Monitoring:} repeat iron studies to confirm response and guide ongoing treatment
    \item \textbf{Dietary advice:} matching intake to needs
    \item \textbf{Education:} on the meaning of results and the plan
\end{itemize}

\section*{Complications}
\begin{itemize}
    \item Unnecessary iron supplementation when deficiency is not confirmed
    \item Missed iron overload with long-term organ damage if untreated
    \item Iron deficiency persisting when the underlying cause is not treated
    \item Misinterpretation when inflammation distorts ferritin
    \item Fatigue and impaired function from untreated deficiency
    \item Complications of unnecessary investigation or treatment
\end{itemize}

\section*{Prognosis and Outcomes}
Iron studies are accurate and reliable when interpreted correctly, and they guide treatment that is highly effective. Iron deficiency corrects with supplementation over weeks to months, and iron overload is managed successfully with regular venesection, preventing the liver, heart, and joint damage of untreated haemochromatosis. Because the tests are cheap, widely available, and informative, they play a central role in the diagnosis and monitoring of iron disorders.

\section*{Prevention and Self-Care}
\begin{itemize}
    \item Have iron studies when you have fatigue, paleness, or unexplained symptoms
    \item Attend for follow-up testing after treatment to confirm correction
    \item Do not take iron supplements on the basis of a single result without medical advice
    \item Report black stools, blood loss, or heavy periods to your doctor
    \item Tell the doctor about inflammation or infection, which affect the results
    \item Keep a record of your results for comparison
    \item Discuss a family history of haemochromatosis with your doctor
    \item Eat an iron-rich diet alongside any treatment
\end{itemize}

\section*{Sources and Further Reading}
The following reputable sources provide further information for healthcare students and patients seeking reliable, current advice about iron studies.
\begin{itemize}
    \item Pathology Tests Explained. \textit{Iron studies.} https://pathologytestsexplained.org.au/
    \item Healthdirect Australia. \textit{Iron studies blood test.} https://www.healthdirect.gov.au/
    \item The Royal Australian College of General Practitioners. \textit{Iron deficiency investigation.} https://www.racgp.org.au/
    \item Haemochromatosis Australia. \textit{Testing and diagnosis.} https://haemochromatosis.org.au/
    \item NHS. \textit{Blood tests for iron.} https://www.nhs.uk/
    \item Lab Tests Online. \textit{Iron tests.} https://labtestsonline.org/
\end{itemize}
